\documentclass{ltjsbook}
\usepackage{docmute} 
\usepackage{../../myStyle} 

\begin{document}
\VerbatimFootnotes
% \chapter{混合分布モデル}
\chapter{Mixed Distribution Model}
\label{x26_mixed}
% さてLM$\to$GLM$\to$GLMM$\to$HLMと，線形モデルを展開させてきました。後半のGLMMやHLMでは分布を混ぜるということをしてきたわけですが，今回も分布を混ぜるモデルについて見ていきたいとおもいます。
So, we have developed the linear model from LM to GLM to GLMM to HLM. In the latter half, with GLMM and HLM, we have been mixing distributions, but this time, I would like to take a look at models that mix distributions.


% 今回の分布の混ぜ方は，パラメータの構造が複雑化するのではなく，そもそも違うパラメータの分布が混ざり合っているというパターンです。具体的な例から見ていきましょう。
The way of mixing in this distribution does not complicate the structure of the parameters, but rather, it is a pattern in which the distributions of different parameters are mixed together. Let's look at it from a specific example.
% \section{混合分布モデル}
\section{Mixed Distribution Model}
% 図\ref{fig:26_01}を見てください。これは野球選手データからTigersの年俸をヒストグラムにしたものです。年俸のデータはそのままだと左に歪んでいるので，対数をとったものを表示しています。
Please refer to Figure \ref{fig:26_01}. This is a histogram of the annual salary of the Tigers made from baseball player data. As the raw salary data is skewed to the left, we are displaying the data after taking its logarithm.
\begin{figure}[ht]
	\begin{center}
		\includegraphics[width=0.8\linewidth,keepaspectratio]{../images/26_mixed/Rplot26_01.png}
		\caption{Tigersの選手年俸(対数)の分布}
		\label{fig:26_01}
	\end{center}
\end{figure}

% ヒストグラムを見ると，どうも2つの山があるように見えますね。赤いラインは密度のカーブを書いたものですが，これはどう見ても2つのピークを持っています。
Looking at the histogram, it seems as though there are two peaks. The red line represents a density curve, and no matter how you look at it, it has two peaks.
% このようなデータに対して，正規分布のモデルを当てはめるのはおかしなことです。\keytermL{正規分布}{せいきぶんぷ}はご存知の通りベルカーブ，すなわちピークは1つ(単峰)で，両裾になだらかに下がっていくものですから，このデータと違う形になっていることは間違いありません。
Applying a normal distribution model to such data is strange. As you know, the normal distribution is a bell curve, thus it has only one peak (unimodal), and it gradually decreases on both sides, so it definitely differs from this data in shape.

% 赤い密度カーブの方が暗示するように，この分布は2つの正規分布が混じり合って出てきているのではないか，と思われます。つまり，出どころや種類の違う分布があるわけです。
As suggested by the red density curve, it appears that this distribution might be derived from a mix of two normal distributions. In other words, there are distributions of different origins or types.
% これまで分析してきたデータは，基本的に同じ分布からデータが出てきたと考えられてきたわけですが，そうでもない分布の場合はどうするか，という問題です。
The data we have analyzed so far has been considered to come from the same distribution, but the question is what to do in the case of a different distribution.

% こうしたモデルを考えるのが\keytermL{混合分布モデル}{こんごうぶんぷもでる}と呼ばれるものです。ここでは複数の正規分布を混ぜ合わせた\keyterm{正規混合モデル}{せいきこんごうもでる}{normal mixture model}を扱います。
The types of models we consider here are referred to as "mixed distribution models". In particular, we will deal with "normal mixture models", which are combinations of several normal distributions.

% データが複数の出自を持つという時に，その違いを示す変数が明確なのであれば困ることはありません。たとえば男性と女性で分布が違うとか，東日本と西日本で分布が違うという既有知識があって，それがデータに含まれていれば，その変数でデータセットを分割して階層モデルにすれば良いだけです。しかし今回のように，阪神の選手という意味で同じ種類のメンバーであるはずなのに，データの表れ方が違う場合，「この2つの分布に分かれる潜在的な変数は何か」ということを探す必要があります。
When the data has multiple origins, it's not a problem if the variables indicating their differences are clear. For example, if there is a known difference in distribution between men and women, or between Eastern and Western Japan, and that information is included in the data, all you need to do is split the dataset by that variable and turn it into a hierarchical model. However, in cases like this, where they should be the same kind of members in terms of being players of the Hanshin team, but the way the data appears is different, there is a need to find out what the latent variables are that separates these two distributions.
% 言い換えれば，データの表面的な特徴からグルーピングをしなければなりません。これはデータの類似性だけを用いて，外的基準を持たずして分類する，\keyterm{クラスタリング}{くらすたりんぐ}{Clustering}という分析方法と関係します。
In other words, you must group the superficial features of the data. This is related to the analysis method called \keyterm{clustering}{clustering}{Clustering}, which classifies without external criteria, using only the similarity of the data.
% そこで少し寄り道になりますが，クラスタリングの手法について見ておきましょう。
So, let's take a brief detour and take a look at the method of clustering.

% \subsection{クラスタリングいろいろ}
\subsection{Various types of clustering}
% クラスタリングという分析方法は，先に述べたように別の変数(外的基準)を使わずに\footnote{\keytermL{機械学習}{きかいがくしゅう}の文脈では，こうした基準となる変数を置かずに分析する手法のことを\keytermL{教師なし学習}{きょうしなしがくしゅう}，と呼びます。}，データの分類を行う方法です。分類は科学の基本で，まず同じものか違うものか，何が同じで何を違うものと考えるか，ということがすべてのスタートになります。クラスタリングではこれをデータから行う必要があり，データの類似性をもとに分類することが基本です。
The method of analysis called clustering, as mentioned earlier, is a method of classifying data without using other variables (external criteria)\footnote{In the context of machine learning, this is called unsupervised learning, a method of analysis without setting these criteria variables.}. Classification is the basis of science, and the starting point for everything is whether things are the same or different, and what is the same and what is considered different. Clustering requires this to be done from the data, and classification based on data similarity is fundamental.

% データの\keytermL{類似度}{るいじど}あるいは同じことですが\textbf{非類似度}を表す指標の基本は\keyterm{距離}{きょり}{distance}です。距離と聞くと普通は2点間を直線で結んだ\keyterm{ユークリッド距離}{ゆーくりっどきょり}{Eucredian distance}を想像されると思いますが，実は距離には他にもいろいろな種類があります(詳しくは第\ref{x15_MDS}講，セクション\ref{distance}，Pp.\pageref{distance}を参照)。たとえば\keytermL{相関係数}{そうかんけいすう}も絶対値を取ればデータの類似度を表す距離情報の一種と考えることができます。この距離を使って，距離が短い＝類似している＝同じクラスター，という考え方から分類していくことになります。
The basis of the index that indicates the similarity or, equivalently, the dissimilarity of the data is distance. When you hear "distance", you might ordinarily imagine the Euclidean distance that connects two points in a straight line, but in fact, there are various other types of distances (For more details, see Lecture \ref{x15_MDS}, Section \ref{distance}, page \pageref{distance}). For instance, if we take the absolute value, correlation coefficient can be considered as a type of distance information expressing data similarity. Using this distance, classifications will be made based on the concept that a shorter distance equals similarity, thus belonging to the same cluster.

% こうした分類，クラスタリングの手法はいくつかあって，大きく分けると\keyterm{階層的クラスタリング}{かいそうてきくらすたりんぐ}{hierarchal clustering}と\keytermL{非階層的クラスタリング}{ひかいそうてきくらすたりんぐ}の二種類に分かれます。階層的というのはHLMの発想と同じで，クラスターが積み重なっていくようなイメージです。たとえば要素$\{a,b,c,d,e\}$をクラスタリングする時，距離が近いものを使ってまず$\{a\},\{b,c\},\{d,e\}$と分割したとします。次に$\{b,c\}$を1つの新しいデータ点だと考えて，$a$と$\{b,c\}$，$\{b,c\}$と$\{d,e\}$の距離を考え，近い方をさらに次のクラスターにまとめます。すなわち$\{a,\{b,c\}\},\{d,e\}$というようにです。それをさらに上位階層でまとめ$\{a,\{b,c\},\{d,e\}\}$とすべてが1つのクラスターにまとめられたら終わり，とより包括的なクラスターへ，より上位のクラスターへとまとめていくのが階層的クラスタリングという手法です。この方法で分類した例を図\ref{fig:26_02}に示します。
There are several methods for classification and clustering, which can be broadly divided into two types: \keyterm{hierarchical clustering}{hierarchical clustering} and \keyterm{non-hierarchical clustering}. The concept of hierarchical is the same as that of HLM, picturing clusters stacking on top of each other. For instance, when clustering elements $\{a,b,c,d,e\}$, suppose we first divide it into $\{a\},\{b,c\},\{d,e\}$ according to the proximity. Then we consider $\{b,c\}$ as a new data point, think about the distances between $a$ and $\{b,c\}$, $\{b,c\}$ and $\{d,e\}$, and group the closer ones into the next cluster. That is, it becomes $\{a,\{b,c\}\},\{d,e\}$. If we continue to group them into higher-level clusters to form $\{a,\{b,c\},\{d,e\}\}$ until all are grouped into one cluster, it ends. This method of grouping into larger clusters, to higher-level clusters, is known as hierarchical clustering. An example of classification using this method is shown in Figure \ref{fig:26_02}.
\begin{figure}[ht]
	\begin{center}
		\includegraphics[width=0.8\linewidth,keepaspectratio]{../images/26_mixed/Rplot26_02.png}
		\caption{Tigersの選手を年俸で階層的クラスタリング(ウォード法による)}
		\label{fig:26_02}
	\end{center}
\end{figure}

% 階層的クラスタリングは，階層が上がる時にクラスター化されたもの同士の距離をどう考えるかによって，いくつかの方法があります。先ほどの例ではまず$\{b,c\}$が1つのクラスターになりましたが，これを$A$とすると，$a$と$A$の距離をどう考えるか，がいろいろあるわけです。大文字の$A$の方は，2つの要素から合成されたもの($A=\{b,c\})$ですから，$a$と$b$の距離を取るのか，$a$と$c$の距離を取るのか，はたまた$b,c$の平均的な値と$a$との距離を取るのか等々，基準が必要なわけです。複数の要素を持つクラスター同士の距離の考え方として，その最大値を取る\textbf{最大距離法}，最小値を取る\textbf{最小距離法}だけでなく，データの重心を取る\textbf{セントロイド法}，平均を取る\textbf{平均法}，クラスター内分散とクラスター間分散の差が最小となるように取る\keyterm{ウォード法}{うぉーどほう}{Ward's method}などがあります\footnote{厳密に書くと，クラスターPに含まれる要素と重心との距離の二乗和を$E(P)$とすると，二つのクラスターP,Qを併合して作られる新しいクラスターTを考えた時，$E(T)-E(P)-E(Q)$をP,Qの距離と考える方法です。}\footnote{Rでは標準関数\verb|hclust|で階層的クラスター分析ができ，ウォード法を使うときは\verb|method='ward.D2'|とします。この\verb|.D2|にはいわれがあって，かつては\verb|ward.D|でウォード法の指定ができたのですが，この関数にバグが含まれていることが判明，修正版を\verb|ward.D2|としたのです。Rは無償のフリーソフトウェアですから，バグが含まれているなんてやっぱり使い物にならない！という批判もあるのですが，オープンなソフトウェアであるからこそソースコードが誰にでも見られて指摘することができ，よくなっていくものでもあります。バグがあったこと，それを自浄作用でフィックスしたことを記録するために，\verb|2|を残すことになりました。プロプライエタリな商用ソフトの場合，間違いがあってもこっそり修正して，同じ関数名で提供されてしまえば，ユーザは分析が間違っていたことに気づくことすらできません。科学技術に関するソフトウェアはオープンであるべきだ，というRの思想がここに現れていると言えるでしょう。}。最後のウォード法は弁別後の結果がわかりやすいため，よく用いられます。
Hierarchical clustering has several methods depending on how we consider the distance between clusters as the hierarchy increases. In the previous example, first $\{b,c\}$ became one cluster, and if we call this $A$, there are various ways to consider the distance between $a$ and $A$. The big $A$ is a composite of two elements ($A=\{b,c\}$), so we need a criterion for whether to take the distance between $a$ and $b$, whether to take the distance between $a$ and $c$, or whether to take the distance between $a$ and the average value of $b, c$, etc. As a concept of the distance between clusters with multiple elements, there is not only the \textbf{maximum distance method} that takes the maximum value and the \textbf{minimum distance method} that takes the minimum value, but also the \textbf{centroid method} that takes the center of gravity of the data, the \textbf{average method} that takes the average, and the \keyterm{Ward's method} that minimizes the difference between the variance within the cluster and the variance between clusters. If you write it rigorously, let $E(P)$ be the sum of squares of the distances between the elements in cluster P and the centroid, and consider P, Q and T as new clusters created by merging the two clusters P, Q. It is a method to consider $E(T)-E(P)-E(Q)$ as the distance of P, Q. In R, you can perform hierarchical cluster analysis with the standard function \verb|hclust|, and when using Ward's method, you specify it as \verb|method='ward.D2'|. This \verb|.D2| has a reason, and once upon a time it was possible to specify Ward's method with \verb|ward.D|, but it was found that this function contained bugs, so the corrected version was \verb|ward.D2|. R is a free software, so there is criticism that it is not useful because it contains bugs, but it is also a software that gets better because it is open and anyone can see the source code and point out. To record that there was a bug and fixed it by self-cleaning, it was decided to leave \verb|2| as is. In the case of proprietary commercial software, even if there is a mistake, if it is quietly corrected and provided with the same function name, the user may not even notice that the analysis was wrong. It can be said that R's idea that software related to science and technology should be open is reflected here. Ward's method is often used because the results after discrimination are easy to understand.

% 一方，階層的でないクラスター分析で有名なのが\keyterm{K-平均法}{K-へいきんほう}{k-means法}です。これは変数間の距離に$k$個の基準点をランダムに置き，その点に近いものをクラスターとします。次に作られたクラスターの重心を新しい基準点として，再度分類します。それの重心を次の基準点とアップデートして再分類$\cdots$と繰り返して，分類結果が変わらなくなる前で何度も反復するのです。だいたい数回の反復で分類できること，データのサイズが大きくても効率よく部類できることからよく利用される方法の1つです。最初に幾つのクラスターに分類するか，明確な基準がないことが欠点でしたが，その後のモデルの展開でクラス数の適合度などを算出して適切なクラス数を求めるようになりました。
On the other hand, the k-means method is a famous non-hierarchical cluster analysis. It places k standard points randomly within the distance between variables and clusters objects that are close to those points. Then, it reclassifies the centroid of the cluster created as a new standard point. It updates the centroid and reclassifies it, repeating this process again and again until the classification result does not change. This method is often used because it can classify in a few iterations and can efficiently divide even large data sizes. The drawback was that there was no clear standard as to how many clusters to classify at first, but later, in the development of models, it started calculating fit parameters such as the number of classes and determining the appropriate number of classes.


% クラスター分析のもう1つの分類基準が，クラスターの境界強度によるものです。\keyterm{ハードクラスタリング}{はーどくらすたりんぐ}{Hard clustering}では，ある個体がどのクラスターに含まれるかがはっきりしています。この個体はクラスターA，この個体はクラスターB，と決められたらそれで決まりです。これに対して\keyterm{ソフトクラスタリング}{そふとくらすたりんぐ}{Soft clustering}は，ある個体がクラスターAに含まれる確率がXX\%，クラスターBに含まれる確率がYY\%，のように分類が確率的で固定的ではない方法です。
Another classification criterion for cluster analysis is based on the strength of cluster boundaries. In hard clustering, it is clear to which cluster a certain individual belongs. Once it has been decided that this individual is in cluster A, and this individual is in cluster B, that's it. Contrary to that, in soft clustering, the classification is not fixed but probabilistic, such as an individual having an XX% chance of being included in cluster A and a YY% chance of being included in cluster B.

% このように，データの特徴から分類をしていくクラスター分析はいろいろな手法があるのですが，今回扱う混合分布モデルはモデルに基づいたクラスタリング，Model based clusteringとも呼ばれる手法です。すなわち，各個体が潜在的な正規分布のどちらに含まれる可能性があるか，その確率を考えながらソフトに分類していく非階層的な方法になります。
In this way, there are various methods of cluster analysis that classify from the features of the data, and the mixed distribution model we will deal with this time is a method called Model based clustering. In other words, it is a non-hierarchical method that classifies softly while considering the probability that each individual may be included in either of the latent normal distributions.

% \subsection{混合分布モデルの考え方}
\subsection{Concept of Mixed Distribution Model}
% 混合分布モデルの考え方は，アルクラスターに含まれるかどうかが確率的に決まるというものです。たとえばコイントスをして，表が出たらクラスターA，裏が出たらクラスターB，というように分類するようなものです。コイントスはベルヌーイ分布に下がいますから，確率$\theta$でクラスターA，確率$1-\theta$でクラスターB，と言っていることと同じです(図\ref{fig:26_03})。
The concept of a mixture distribution model is that whether it is included in a cluster is probabilistically determined. For example, flip a coin, and if you get heads, it's Cluster A, and if you get tails, it's Cluster B, that's how classification is done. Coin tossing is under the Bernoulli distribution, so it's the same as saying that it's Cluster A with probability $\theta$, and Cluster B with probability $1-\theta$ (Figure \ref{fig:26_03}).
\begin{figure}[ht]
   \begin{center}
      \includegraphics[width=0.8\linewidth,keepaspectratio]{../images/26_mixed/slide26_01.png}
      \caption{二つの正規分布からデータが出てくるモデル}
      \label{fig:26_03}
   \end{center}
\end{figure}


% もちろん3つ以上のクラスターに分類することもあるでしょう。この場合は\keytermL{カテゴリカル分布}{かてごりかるぶんぷ}(セクション\ref{DDist}，Pp.\pageref{DDist}を参照)に従う確率変数$z$を考えることになります。
Of course, there may be cases where we classify into more than three clusters. In this case, we would consider a probability variable $z$ that follows a categorical distribution (see section \ref{DDist}, Pp.\pageref{DDist}).

% データがK個の正規分布，いずれかから得られていると考えるとしましょう。ここで各正規分布はそれぞれの位置母数$\mu_k$とスケール母数$\sigma_k$を持っているとします。個々のデータ$i$がどの正規分布からきているかを表す混合確率$\lambda_{ik}$を考えると，$\lambda_{ik} \ge0$で$\sum_{k=1}^K \lambda_{ik} = 1$になります。つまり，$i$ごとのK-simplexベクトルを考えることになります。
Let's assume that the data is obtained from K normal distributions. Here, each normal distribution has its own location parameter $\mu_k$ and scale parameter $\sigma_k$. If we consider the mixture probability $\lambda_{ik}$, which indicates which normal distribution each data $i$ comes from, then $\lambda_{ik} \ge0$ and $\sum_{k=1}^K \lambda_{ik} = 1$. In other words, we have to consider a K-simplex vector for each $i$.
% 個々のデータ$Y_i$の出自は，$z[i] \sim Categorical(\lambda_{i})$で決まり，データ$Y_i \sim N(\mu_{z[i]},\sigma_{z[i]})$からきている，と考えることになります(図\ref{fig:26_04})。
The origin of each data point $Y_i$ is determined by $z[i] \sim Categorical(\lambda_{i})$, and it comes from the data $Y_i \sim N(\mu_{z[i]},\sigma_{z[i]})$ (Figure \ref{fig:26_04}).

\begin{figure}[ht]
	\begin{center}
		\includegraphics[width=0.8\linewidth,keepaspectratio]{../images/26_mixed/slide26_02.png}
		\caption{K個の正規分布からデータが出てくるモデル}
		\label{fig:26_04}
	\end{center}
\end{figure}

% データ生成メカニズムの設計図は図\ref{fig:26_03}，\ref{fig:26_04}の通りです。具体的なプログラミングに進む前に，イメージをしっかり掴んでおきましょう。
The design of the data generation mechanism is as shown in figures \ref{fig:26_03} and \ref{fig:26_04}. Make sure to firmly grasp the image before proceeding with specific programming.
% これらの設計図の一番上にある，どのクラスターからデータが出てくるかを確率的に決めるというところがポイントです。ですがこのモデルを実装できるようになれば，データ発生の条件分岐を確率モデルで表現できるようになるのです。
The key point is at the top of these blueprints, where it's decided probabilistically from which cluster the data comes out. But once you can implement this model, you can express conditional branching of data generation with a probability model.
% 確率$\theta$で左，$1-\theta$で右のルートを取る，というように，データが出てくるメカニズムの背後に潜在的な条件を設定できますから，分析の可能性がグッと広がることになります。
You can assume underlying conditions behind the mechanism of data generation, such that a route to the left is taken with a probability of $\theta$, while a one to the right is taken with a probability of $1-\theta$. This will significantly expand the possibilities of analysis.
% 実践的にはたとえばマーケティングにおいて，顧客がどういうクラスター(主婦層，独身者，大家族etc...)に属するかはわかりませんが，買い物の傾向から分類するというようなことができます。あるいは心理学的なシーンでも，潜在的に異なる種類(抑うつ傾向があるタイプ，ないタイプなど)が，反応傾向の違いから読み取れるようになるかもしれません。
In practical terms, for example in marketing, you may not know which cluster (housewives, singles, large families etc...) a customer belongs to, but you can categorize them based on their shopping trends. Alternatively, in psychological scenarios, it might be possible to delineate different types (such as those with a tendency towards depression, those without) based on differences in responsiveness.

% クラスター分析はデータから分類するので，分類後にそれがどう言った変数と関係しているのか，あるいは分類されたクラスターにはどういう特徴があるのかを探索することが一般的です。それでもデータだけから分類できるのは魅力的な手法であると言っていいでしょう。
Cluster analysis is typically used to explore variables that may relate to data classification, or the characteristics of classified clusters, as it classifies data. Still, it can be said that it is an attractive method as it classifies only from the data.

% \section{ターゲット記法と周辺化消去}
\section{Target Notation and Peripheral Elimination}
% さてこの条件分岐をどのようにコードに書いていくか，というところに話を進めましょう。ここで新しく，2つの技術を導入する必要があります。1つはターゲット記法，もう1つは\keyterm{周辺化消去}{しゅうへんかしょうきょ}{marginalizing}です。
Now, let's move on to how to write this conditional branching in code. Here, we need to introduce two new techniques. One is target notation, and the other is marginalization.
% \subsection{ターゲット記法}
\subsection{Target Notation}
% これまでStanで，確率変数に従う分布の記法は$\sim$を使って表現していました。すなわち，コード\ref{stan::26_01}のような書き方です。
Up until now, the notation for distributions following random variables in Stan has been expressed using $\sim$, as in the writing method in code \ref{stan::26_01}.
\begin{lstlisting}[language=C++,label={stan::26_01},caption={sampling記法}]
for(i in 1:N){
    Y[i] ~ normal(mu,sigma);
}
\end{lstlisting}

% これと同じ動きをする，別の表記方法があります。それが\keytermL{ターゲット記法}{たーげっときほう}というもので，コード\ref{stan::26_02}のように書きます。
There is another notation that works the same way. This is called the "target notation", and it is written as in code \ref{stan::26_02}.
\begin{lstlisting}[language=C++,label={stan::26_02},caption={target記法}]
for(i in 1:N){
    target += normal_lpdf( Y[i] | mu, sigma);
}
\end{lstlisting}


% 何やら奇妙な書き方に見えますが，順番に見て行きましょう。まず\verb|+=|の右側，\verb|noraml|が\verb|norma_lpdf|になっているところに注目してください。ここでlpdfはlog probability density function，すなわち対数確率密度関数という意味です。\verb|normal|というのがついていますから，正規分布の対数確率密度関数ということになりますね。ちなみに離散変数の場合は\keyterm{確率密度}{かくりつみつど}{probability density}ではなく\keyterm{確率質量}{かくりつしつりょう}{probability mass}になりますから，lpmf，すなわちlog probability mass function(対数確率質量関数)になります。具体的には，\verb|bernoulli_lpmf|とか\verb|poisson_lpmf|となりますので注意してください。
This looks like a strange way of writing, but let's look at it in order. First, pay attention to the right side of \verb|+=|, where \verb|noraml| has become \verb|norma_lpdf|. Here, lpdf means log probability density function, that is, a logarithmic probability density function. Since \verb|normal| is attached, it means a normal distribution logarithmic probability density function. By the way, in the case of discrete variables, it is not \keyterm{probability density} but \keyterm{probability mass}, so it becomes lpmf, that is, log probability mass function. Specifically, it will be \verb|bernoulli_lpmf| or \verb|poisson_lpmf|, so please be careful.

% さて対数確率密度関数ってのがなぜ出てきたのでしょうか。ここで\keyterm{尤度}{ゆうど}{likelihood}のことを思い出してほしいのですが，尤度とはデータが\keytermL{確率分布}{かくりつぶんぷ}から出てくるときの尤もらしさ，出てきやすさのような指標なのでした。確率分布はパラメータがわかっている時にデータが出てくる確率を記述する関数ですが，同じ関数でパラメータが未知，データが既知の場合(多くの研究シーンはこちらですが)に，それを尤度関数とよぶのでした。データが得られた時の尤度を計算し，その尤度が最も大きくなるところをパラメータの推定値とするのが\keyterm{最尤推定法}{さいゆうすいていほう}{Maximum likelihood method}であり，尤度を使って事後分布を算出，事後分布の形から推定値を考えるのが\keyterm{ベイズ推定法}{べいずすいていほう}{Bayesian inference method}なのでしたね。つまりMCMCで計算するにはデータ全体の尤度を考えているのです。そしてデータ全体尤度は各データ点の尤度の総積$\prod$で考える必要がありますが，確率の掛け算は数字が小さくなるので\keyterm{対数尤度}{たいすうゆうど}{log-likelihood}の総和で計算するのでした。
So why did the log-likelihood density function come in? Please remember what the \keyterm{likelihood} was, which is an indicator of how likely data comes from the \keytermL{probability distribution}, that is, how likely it appears. The probability distribution is a function that describes the probability that data comes out when parameters are known, but the same function is called the likelihood function when parameters are unknown and data is known (most of the research scenes are like this). When data is obtained, calculate the likelihood and use the \keyterm{maximum likelihood method}, which is where the likelihood is maximized as the estimate of the parameter, and use the likelihood to calculate the posterior distribution, and consider the estimated value from the shape of the posterior distribution This was the \keyterm{Bayesian inference method}. In other words, MCMC is considering the likelihood of all the data. The whole data likelihood needs to be considered as the total product $\prod$ of the likelihoods of each data point, but because the multiplication of probabilities results in small numbers, it was calculated as the total sum of \keyterm{log-likelihood}.

% 具体的に書くと，データ$\bm{Y}=Y_1,Y_2,Y_3$がパラメータ$\theta$のベルヌーイ分布から出てきている場合，データの尤度は
To be specific, when the data $\bm{Y}=Y_1,Y_2,Y_3$ is coming from a Bernoulli distribution with parameter $\theta$, the likelihood of the data is
\[ L(\bm{Y}|\theta) = \prod_{i=1}^3 bernoulli(Y_i|\theta) = bernoulli(Y_1|\theta) \times bernoulli(Y_2|\theta ) \times bernoulli(Y_3|\theta) \]
であり，対数尤度は
\[ LL(\bm{Y}|\theta) = \sum_{i=1}^3 \log(bernoulli(Y_i|\theta)) = \log(bernoulli(Y_1|\theta))+\log(bernoulli(Y_2|\theta))+\log(bernoulli(Y_3|\theta)) \]
% となるわけです。ここで$bernoulli(Y|\theta)$の縦棒$\mid$は，$\theta$が与えられた時の$Y$の出る確率，という条件付き確率の記号になります。
So it turns out. Here, the vertical bar $| in $bernoulli(Y|\theta)$ represents the conditional probability symbol of the occurrence of $Y$ when $\theta$ is given.

	% プログラミング的には，$Y_1,Y_2,\cdots,Y_n$は\verb|for(i in 1:N){ ~ }|という形で1からNまでの繰り返し計算を意味しますから，コード\ref{stan::26_01}や\ref{stan::26_02}でやっていることは，各データ点についての対数尤度を足し合わせて行っている作業そのものを意味しています。
	In terms of programming, $Y_1,Y_2,\cdots,Y_n$ means repeating calculations from 1 to N in the form of \verb|for(i in 1:N){ ~ }|. So, what we are doing in code \ref{stan::26_01} and \ref{stan::26_02} is summing the log likelihoods for each data point, which is exactly what we are doing.

	% ここでターゲット記法にある\verb|+=|ですが，これはプログラミング独特の表記方法で，数学表記ではありません。プログラミングでは代入を使って，たとえば変数Aに1を加えたものを新しいAとする(上書きする)，ということを\verb|A = A + 1|と表記できます。この\verb|+=|はこうした上書きを一言で書く記号で，このAの例だと\verb|A += 1|とすれば1を加えて上書き，と同じ意味になります。
	Here we have \verb|+=| in the target notation, which is a unique way of representing in programming, not a mathematical notation. In programming, you can use assignment to express, for example, adding one to a variable A and making it the new A (overwriting it) as \verb|A = A + 1|. This \verb|+=| is a symbol that expresses such overwriting, and in the case of this A example, \verb|A += 1| means adding one and overwriting, which has the same meaning.
	% これを踏まえてターゲット記法を見ると，\verb#target += normal_lpdf( Y[i] | mu, sigma)#ですから，\verb#target = target + normal_lpdf( Y[i] | mu, sigma)#としていることと同じです。つまり，対数尤度を足しあげて行ってね，ということを明示的に書いていることになります。\verb|target|は足し合わされたもの，というだけで「目標はこれ」を意味するぐらいの予約語だと思ってください。
	Considering this and looking at the target notation, \verb#target += normal_lpdf( Y[i] | mu, sigma)#, it's the same as writing \verb#target = target + normal_lpdf( Y[i] | mu, sigma)#. In other words, it explicitly writes that you should add up the log-likelihood. Please think of \verb|target| as a reserved word that means 'the summed up ones', or something like 'this is the goal'.

	% サンプリング記法とターゲット記法は，書き方が違えど同じ働きをするのですからどちらでもいいじゃないか，と思うかもしれませんが，ターゲット記法でないと表現できないようなことがあるので，このような代替手法が用意されているのだと思ってください。ターゲット記法でないと表現できないようなこと，というのはもちろん今回の，分布を混ぜ合わせる時がそれです\footnote{そのほかにもターゲット記法は，サンプリング記法では省略されている事後対数確率も記録するという特徴があり，ブリッジサンプリングによる事後対数尤度の計算などにはこちらの表記が必要になってきます。もっとも非常に専門的なことなので，ここでは条件分岐の際に必要な記法，と割り切っていただいても構いません。}。
	You might think that the sampling notation and the target notation, which do different jobs despite their different notations, are equally good. However, there are certain things that can only be expressed by target notation, so please consider this as an alternative method. The case where target notation is necessary for expressing things, of course, refers to the current situation of mixing distributions\footnote{In addition, the target notation also has a feature of recording posterior log probabilities that are omitted in the sampling notation, which becomes necessary for calculations such as posterior log likelihood by bridge sampling. This is a very specialized topic, so you may just consider it as a notation needed for conditional branching at this point.}.
	% \subsection{周辺化消去}
	\subsection{Elimination of Peripheralization}
	% さてここまでさまざまな問題を解決してくれたStanですが，Stanには離散型のパラメータを扱うことができない，という欠点があります\footnote{JAGSという確率的プログラミング言語でしたら，これは可能です。}。
	Now, Stan has helped us solve various problems up to this point, but it has a disadvantage that it can't handle discrete parameters\footnote{This is possible with a probabilistic programming language called JAGS.}.
	% 今回は$z[i] \sim Categorical(\lambda_{i})$と，どの目が出たかというパラメータが離散的ですから，このままではコード化できません。ではどうするか。仕方がないので，「起こりうる可能性をすべて数え上げて考える」という方法を取ることにします。
	This time, since the parameter $z[i] \sim Categorical(\lambda_{i})$ indicating which outcome occurred is discrete, it cannot be coded as it is. So what do we do? There's no other way but to adopt a method of "enumerating and considering all possible outcomes".
	% 最終的な事後分布は対数尤度関数を足し合わせたものになりますから，条件分岐したルートすべてについて対数尤度関数を考えて，それを足し合わせてやれば良いのです。
	The final posterior distribution is the sum of the log-likelihood functions, so all you have to do is consider the log-likelihood function for all the conditioned branched routes and add them up.

	% 簡単な例として図\ref{fig:26_03}にあるような，2つの正規分布モデルを考えてみましょう。コイントスで表が出れば$N(Y_i \mid \mu_1,\sigma_1)$，裏が出れば$N(Y_i \mid \mu_2,\sigma_2)$のルートに行けば良いのです。表が出る確率は$\theta$ですから，表が出るルートを通ってデータ$Y_i$が出てくる確率は$\theta \times N(Y_i\mu_1,\sigma_1)$になります。同じく裏が出るルートを通ってデータが出てくる確率は，$(1-\theta) \times N(Y_i \mid \mu_2,\sigma_2)$ということになります。
	Let's consider two normal distribution models as a simple example, as shown in Figure \ref{fig:26_03}. If a coin flip results in heads, you should follow the $N(Y_i\mid\mu_1,\sigma_1)$ path, and if it results in tails, you should follow the $N(Y_i\mid\mu_2,\sigma_2)$ path. The probability of getting heads is $\theta$, so the probability of getting data $Y_i$ by following the path when heads is obtained is $\theta \times N(Y_i\mid\mu_1,\sigma_1)$. Similarly, the probability of getting data by following the path when tails are obtained is $(1-\theta)\times N(Y_i\mid\mu_2,\sigma_2)$.


	% ですから，最終的にデータ$Y_i$が出てくる確率は$\theta \times N(Y_i \mid \mu_1,\sigma_1) + (1-\theta) \times N(Y_i\mid\mu_2,\sigma_2)$だ，ということができるわけです。
	Therefore, the probability of ultimately getting the data $Y_i$ can be expressed as $\theta \times N(Y_i \mid \mu_1,\sigma_1) + (1-\theta) \times N(Y_i\mid\mu_2,\sigma_2)$.

	% さて，これらは確率をそのまま使った表現になっていますが，Stanの中では対数で考えるのでした。つまり表が出るルートの対数尤度は
	Now, these are expressions that use probability as it is, but in Stan, we think in terms of logarithms. In other words, the log-likelihood of the route where the table appears is
	\[ \log(\theta) + \log(N(Y_i\mid\mu_1,\sigma_1)) \]
	ですし，裏が出るルートもまた同様に$\log$をとって考える必要があります。対数を取ると積が和になるので，$\times \to +$に変わっていることに注意してください。

	ところがこの表ルート，裏ルートは足し合わせないといけません。対数をとったものを直接足し合わせるわけにはいきませんから，対数尤度をただの尤度に戻す，つまり$\log$の逆関数である$\exp$関数を通す必要があります。ややこしいですが，
	\[ \exp\left( \log(\theta) + \log(N(Y_i\mid\mu_1,\sigma_1)) \right) + \exp\left( \log(1-\theta) + \log(N(Y_i\mid\mu_2,\sigma_2)) \right) \]
	% が尤度であり，これの対数をとったものを\verb|target|に加える必要があるのです。
	This is the likelihood, and you need to add the logarithm of this to \verb|target|.

	% このように，$\exp$を通して，足し合わせ($\sum$)て，さらにその$\log$をとるという作業をひとまとめにした関数が，Stanにはあります。それが\verb|log_sum_exp|関数です。
	Thus, through $\exp$, a function that encompasses the operations of addition ($\sum$) and further taking its $\log$, we have this function in Stan. This function is the \verb|log_sum_exp| function.
	% この関数は，\verb|log_sum_exp(x,y)|とすると$\log(\exp(x)+\exp(y))$という操作をします。この関数に渡すものがベクトルであれば，
	This function, \verb|log_sum_exp(x,y)|, performs an operation called $\log(\exp(x)+\exp(y))$. If the values passed to this function are vectors,

	\[ \verb|log_sum_exp|\left( \begin{pmatrix} x_1\\x_2\\ \vdots \\ x_n \end{pmatrix} \right) = \log\left(\sum_{i=1}^n exp(x_i) \right) \]

	という操作をすることになります。

	これを使って混合分布モデルの実装をしていきましょう。
	\subsection{混合分布モデルのコード}
	先ほどの，阪神タイガースにおける選手の年俸分布を2つの正規分布から出てきているもの，と考えてコード化したものは，Code\ref{stan::26_03}のようになります。
	\begin{lstlisting}[language=C++,label={stan::26_03},caption={混合正規モデル}]
data {
  int<lower=1> K;
  int<lower=1> L;
  array[L] real Y;
}

parameters {
  array[L] simplex[K] theta;
  ordered[K] mu;
  array[K] real<lower=0> sigma;
}

transformed parameters{
  array[L] vector[K] lp;
  for (l in 1:L) {
    for (k in 1:K) {
      lp[l,k] = log(theta[l,k])+ normal_lpdf(Y[l]|mu[k],sigma[k]);
    }
  }
}

model{
  for(l in 1:L){
    target+=log_sum_exp(lp[l]);
  }
  sigma ~ cauchy(0,5);
  mu ~ normal(0,10);
}

generated quantities{
  array[L] vector[K] prob_class;
  array[L] int<lower=1,upper=K> pred_class;
  for(l in 1:L){
    prob_class[l] = softmax(lp[l]);
    pred_class[l] = categorical_rng(prob_class[l]);
  }
}
\end{lstlisting}
	\paragraph{コード解説}
	\begin{description}
		\item[dataブロック] クラスター数K，データ長L，アウトカム変数Yを宣言しています。
		\item[parametersブロック] 合計が1になるKの長さを持つベクトルである，\verb|simplex|型ベクトルを，データ数だけ用意します。また，K個の正規分布からデータが出てきているのですから，平均と標準偏差もK個用意します。ここで大事なポイントとして，平均$\mu$が\verb|ordered|型としてあるところです。このベクトルは，要素$\mu_1>\mu_2>\cdots>\mu_k$と，要素が大きさの順に並んでいることです。このような制約をかけておかないと，K個の正規分布がどの位置にあるのか判然としな苦なるという問題が生じます。たとえばあるチェインで$\mu_1>\mu_2$となっていて，データ$X_k$が1番目の正規分布から出てきていると考えられたとしましょう。しかし別のチェインで$\mu_1<\mu_2$となっていれば，$X_k$は平均値が大きい方から出てきているのですから，こちらのチェインでは2番目の正規分布から出てきたものとして，MCMCサンプルが進むことになります。複数のチェインから得られたMCMCサンプルは最終的に統合されますから，このようになっていると$X_k$は正規分布1からも2からも同じように出てきていることになります。各チェインでは正しくサンプリングできているのですが，チェインごとに分布のラベルが違うので，統合すると訳がわからなくなる，この問題をとくに\keyterm{ラベルスイッチング}{らべるすいっちんぐ}{Label Switching}問題といい，混合分布モデルではよく生じる問題なのです。そこで分布の位置母数に順番の制約を入れて，ラベルの変動を止める必要があります。
		\item[transformed parametersブロック] 長さKのベクトルをデータの数Lだけ用意した\verb|lp|変数を作りました。これはlog-probabilityの略でこのような名前にしましたが，任意の名前で結構です。l行目のデータがクラスターkに含まれる可能性は，$\theta_{lk} \times N(Y[l]\mid\mu_k,\sigma_k)$ですから，それの対数をとった形で表現しています。
		\item[modelブロック] モデル尤度のところは，対数$\to$ $\exp$関数で確率の形 $\to$ ベクトルの要素をすべて足し合わせる $\to$ $\log$をとって対数尤度にする，という計算をまとめて\verb|log_sum_exp|関数で行い，それをすべて\verb|target|に追加していくという方法で書いています。事前分布のところは普通のサンプリング記法で書きました。ちなみに$\mu,\sigma$はそれぞれK個ずつありますが，Stanは添字の省略を許してくれます。
		\item[generated quantitiesブロック] 久しぶりに生成量ブロックの登場です。ここでは分析結果を使って，l行目のデータが結局何番目の正規分布から出てきたのかを逆算して出しています。まず，対数尤度で表現されているK個の要素を持つベクトルを，確率の形に計算したものを\verb|preb_class|変数にしています。ここで使われている\verb|softmax|という関数は，
		      \[ softmax(x) = \frac{exp(x)}{\sum_{k=1}^K exp(x_k)}\]
		      %       という計算をするものです。分母はベクトルの総和で分子がその要素ですから，要するに各クラスに入る確率を表すベクトルになっている訳です。これを使って\verb|pred_class|変数を作っています。今回のデータから予想される所属クラスなので，K面サイコロを振った出目という形で表現します。\verb|categorical_rng|はカテゴリカル分布からの乱数発生で，その発生確率が先ほどの\verb|prob_class|になっているので，この乱数で出た目が予測されるクラスということです。
		      This performs calculations. The denominator is the sum of the vectors and the numerator is its element, which essentially a vector representing the probability of entering each class. Using this, we create a variable called \verb|pred_class|. As it is a predict class based on the current data, it is expressed as the outcome of rolling a K-faced die. \verb|categorical_rng| is a random number generation from a categorical distribution, and since the generation probability is the above mentioned \verb|prob_class|, the predicted class is determined by this random number.
	\end{description}
	% これを実行するためのRコードをコード\ref{code:26_01}に用意しました。
	I have prepared R code in Code\ref{code:26_01} to execute this.
	\begin{lstlisting}[language=c++,label={code:26_01},caption={混合分布モデルのコード}]
dat <- read_csv("baseballDecade.csv")
dat.tmp <- dat %>%
  filter(position != "投手") %>%
  filter(Games > 50) %>%
  filter(team == "Tigers") %>%
  mutate(salary.log = log(salary)) %>%
  dplyr::select(salary.log, Name) %>%
  rowid_to_column("ID")

dataSet <- list(K = 2, L = NROW(dat.tmp), Y = dat.tmp$salary.log)
model <- cmdstanr::cmdstan_model("latent.stan")
fit <- model$sample(data = dataSet, chains = 4, parallel_chains = 4)
\end{lstlisting}


	% \paragraph{コード解説}
	\paragraph{Code Explanation}
	\begin{description}
   \item[1行目] データファイル\verb|baseballDecade.csv|を読み込みます。
   \item[2-8行目] 分析用のデータに加工しています。
      \begin{description}
         \item[3行目] フィルターをかけてデータを野手だけのものにします。
         \item[4行目] フィルターをさらにかけて，50試合場出ている選手に限定します。
         \item[5行目] チームがTigersの選手だけに絞っています。
         \item[6行目] 年俸のデータは歪んでいるので，対数をとって正規分布の形に近似させました\footnote{データを変形させずに確率分布の方で対応するのがよい，と常々話をしてきましたが，ここではデータを変形させています。実は左に歪んだような分布は対数正規分布というのがあり，そこからのデータ生成を考えても良いのですが，対数正規分布はデータの対数が従う正規分布というだけなので，データの方を変えてしまった方が説明しやすいのです。比率や度数など，データの生成メカニズムがそもそも違うばあいは，データを変換せずに使うのが正しいやり方です。}。
         \item[7行目] 使う変数だけに絞っています。
         \item[8行目] 行番号を\verb|ID|という変数名にして追加しました
      \end{description}

	%    \item[10行目] データセットを作っています。
	\item[10th line] We are creating a dataset.

	%    \item[11行目] cmdstanでコンパイルしています。rstanパッケージを使う場合はrstanのコンパイル関数を使ってください。
	\item[Line 11] Compiling with cmdstan. If you are using the rstan package, please use the compilation function of rstan.
	%    \item[12行目] cmdstanのサンプリング関数です。rstanパッケージを使う場合はrstanのサンプリング関数を使ってください。
	\item[Line 12] This is a sampling function of cmdstan. If you use the rstan package, please use the sampling function of rstan.
	\end{description}

	% 推定した結果を見てみましょう。私の環境では出力\ref{MCMCout:out26_01}のようになりました。
	Let's take a look at the estimated results. In my environment, it turned out as shown in output \ref{MCMCout:out26_01}.
	\begin{MCMCout}{}{out26_01}
	\begin{verbatim}
# A tibble: 4 × 7
  name           EAP       MED       MAP        SD       L95       U95
  <chr>    <num:.3!> <num:.3!> <num:.3!> <num:.3!> <num:.3!> <num:.3!>
1 mu[1]        7.624     7.624     7.638     0.107     7.415     7.826
2 mu[2]        9.702     9.699    10.017     0.284     9.164    10.138
3 sigma[1]     0.656     0.653     0.646     0.063     0.544     0.788
4 sigma[2]     0.762     0.786     0.901     0.237     0.380     1.192
\end{verbatim}
\end{MCMCout}

	% 対数をとった時の平均値が$7.624$と$9.702$ですから，指数関数を入れて読み直せば$\exp(7.624)=2046.733$,$\exp(9.702)=16350.28x$となります。2000万円クラスの選手と，1億6500万円クラスの選手の2つのグループに分かれることが見て取れますね。
	The average values when taking the logarithm are $7.624$ and $9.702$, so if you re-read it with an exponential function, it becomes $\exp(7.624)=2046.733$, $\exp(9.702)=16350.28x$. We can see that it is divided into two groups, players in the 20 million yen class and players in the 165 million yen class.
	% 事後予測分布から，所得平均の高い「超一流」選手と「一般的」な選手とに分けて分布を色分けしてみました(図\ref{fig:26_10})。データから，この選手がどちらのクラスに所属しているかも想像できるところがおもしろいですね。
	I have color-coded the distribution by separating the players with high average incomes, the 'super elite', from the 'typical' players, based on the post-prediction distribution (Figure \ref{fig:26_10}). It's interesting that we can guess which class this player belongs to from the data.
	\begin{figure}[ht]
	\begin{center}
		\includegraphics[width=0.8\linewidth,keepaspectratio]{../images/26_mixed/Rplot26_06.png}
		\caption{選手のクラスタリング}
		\label{fig:26_10}
	\end{center}
\end{figure}




	% \section{ゼロ過剰ポアソン分布モデル}
	\section{Zero-Inflated Poisson Distribution Model}
	\label{ZIpoisson}
	% さて，異なる分布の混ぜ合わせ，離散確率分布による条件分岐の例として，もう1つ別の例を示しましょう。次の図\ref{fig:26_05}を見てください。
	Now, let's show another example of a mixture of different distributions and conditional branching by a discrete probability distribution. Please refer to the following Figure \ref{fig:26_05}.
	\begin{figure}[hb]
	\begin{center}
		\includegraphics[width=0.8\linewidth,keepaspectratio]{../images/26_mixed/Rplot26_03.png}
		\caption{投手のセーブ数}
		\label{fig:26_05}
	\end{center}
\end{figure}


	% これは投手のセーブ数の度数分布です。野球は9回のターン制で，1回の最初から投げる投手は先発投手と言います。先発投手が9回まで投げ切れればいいのですが，疲労も溜まるし戦術によっては投手の打席に代打を投入することもあります。投手が降板した場合は，次の投手が代わって続きから投げますが，「勝っている試合で後半まできたので，最後の1，2回は確実に勝ち逃げできるようにしたい」というときに投入されるのが抑え(クローザー)と呼ばれる投手です。最近のプロ野球は分業が進んでいますから，先発投手が完投することは少なく，100球程度で交代し，中継ぎ--抑えと投手リレーしていくのが一般的な戦術です。さて，先発投手は最初から出てきて6，7回まで投げますから(途中でボロボロに打たれたり怪我で降板などがなければ，ですが)，毎日のように登板することはなく，一度登板すると3，4日休んで次の試合にでる，ということになります。日本プロ野球は年間140試合近くありますが，通年で2-30回も登板することはありません。一方，抑えの投手は最後の8，9回を投げる程度で，しかも負けている試合などでは出番がありません。必然的に中4日の休憩なども入りませんから，試合数はたくさん出ることができます。そして抑えに成功するとセーブポイントというのがつきますが，このセーブポイントは優秀なクローザーであれば年間数十ポイント取ることができる訳です。
	This is a frequency distribution of the number of saves by a pitcher. Baseball is played in nine rounds, and the pitcher who throws from the beginning of the first round is known as the starting pitcher. It would be great if the starting pitcher could pitch until the ninth round, but fatigue builds up and there are times when a pinch hitter is brought in for the pitcher's slot for tactics. When the pitcher is relieved, the next pitcher takes over from there, but when it's a case of; "We're winning the game so far, so we want to make sure we can safely pull away in the last one or two innings", that's when a pitcher called a closer is brought in. Pro baseball today is highly specialized, so it's less common for starting pitchers to throw complete games and they usually switch after about 100 pitches, and the typical strategy is a chain of starting pitcher - relief pitcher - closer. Now, the starting pitcher, who starts from the beginning and throws until the 6th or 7th inning (assuming they're not hit hard or injured midway), won't pitch every day, and once they pitch, they'll take a 3 or 4 day rest before the next game. There are nearly 140 games in a season of Japanese professional baseball, but it's unusual to pitch even 20-30 times a year. On the other hand, closers throw the last 8th or 9th inning, and they don't even get a turn in a losing game. They won't necessarily get a four-day rest, so they can play many games. When they succeed as a closer, they earn what's called save points, and that means a skilled closer can earn dozens of points a year.

	% さて，図\ref{fig:26_05}にあるようにこのセーブポイントを見てみますと，ほとんどの投手がゼロになっています。以下データをみますと，1セーブが53人，2セーブが26人・・・となって，20〜40セーブある投手が10人弱ずつぐらいいる，という分布をしています。これは先ほども言った分業制のせいで，先発投手にセーブポイントがつくことはありませんし，先発投手は1チームに10人弱ぐらいは揃えているでしょうから，投手全体で見るとセーブポイントは0の人が多いのです。
	Now, when we look at this save point shown in Figure \ref{fig:26_05}, almost all of the pitchers are zero. Looking at the data below, we see that there are 53 people with one save, 26 people with two saves... and about 10 people with 20 to 40 saves, showing this distribution. This is due to the specialization system I mentioned earlier, so starting pitchers don't get save points, and since each team will have about ten starting pitchers, most pitchers will have zero save points when we look at them as a whole.

	% このセーブポイント，0以上の整数を取りますからポアソン分布に従うと考えられるのですが，ポアソン分布のパラメータ$\lambda$をどう変えても(図\ref{fig:26_06})，こんな形にはなりません。
	This save point can take any non-negative integer, so it can be thought to follow a Poisson distribution. However, no matter how we change the parameter $\lambda$ of the Poisson distribution (figure \ref{fig:26_06}), it does not take this shape.


	\begin{figure}[ht]
	\begin{center}
		\includegraphics[width=0.8\linewidth,keepaspectratio]{../images/26_mixed/Rplot26_04.png}
		\caption{さまざまなポアソン分布。$\lambda$は1.5，3，4.5，10.5で描画}
		\label{fig:26_06}
	\end{center}
\end{figure}

	% そこで混合分布です。まずセーブ数が0かどうかを考えます。セーブ数0というのは，そもそもクローザーではないか，クローザーなのに0セーブポイントのへっぽこか，の二択です。セーブ数が0でない限り，ポアソン分布からセーブ数が生成される，と考えるのです。設計図のイメージは図\ref{fig:26_07}の通りです。
	So, this is a mixed distribution. First, let's consider whether the save number is zero or not. Having a save number of zero means either the player is not a closer at all, or they are a closer but a total dud with zero save points, it boils down to these two possibilities. As long as the save number is not zero, it is generated from a Poisson distribution, that's our assumption. The image of the blueprint is as shown in Figure \ref{fig:26_07}.
	\begin{figure}[ht]
	\begin{center}
		\includegraphics[width=0.8\linewidth,keepaspectratio]{../images/26_mixed/slide26_03.png}
		\caption{ポアソン分布との混ぜ合わせ}
		\label{fig:26_07}
	\end{center}
\end{figure}


	% このモデルがちょっと特殊なのは，ご覧の通り0というデータの値一箇所に極端な偏りがあるところです。この分布のことを\keyterm{ゼロ過剰ポアソン}{ぜろかじょう}{Zero-Inflated Poisson}といいます。
	The unique feature of this model, as you can see, is the extreme bias at a data value of 0. We refer to this distribution as the "Zero-Inflated Poisson".

	% 注意してほしいのですが，ポアソン分布から0の度数が出てくることもあり得ます。逆に考えると，データ$Y_i=0$であれば，これはクローザーでないからそうなっているのか，クローザーなのにヘッポコでセーブ数がつかなかったのかのどちらか，という判断をしなければならないということです。つまり，データからの分岐を考えることになります。クローザーかどうかは確率$\theta$で決まるとして，$\theta$ならクローザー，$1-\theta$ならクローザーじゃないとすると，
	Please note that a frequency of 0 can come out from a Poisson distribution. Conversely, if the data $Y_i=0$, it means we have to decide whether it's because they're not a closer or they're a closer but were not good enough to improve the save total. In other words, we need to think about branching from the data. If it's determined by the probability $\theta$ whether it's a closer or not, let's say it's a closer for $\theta$ and not a closer for $1-\theta$.
\begin{equation}
	\left\{ \begin{array}{cc} \mbox{データが0ではない} & \to \theta \times Poisson(Y_i\mid\lambda)                    \\
             \mbox{データが0だ}         & \to 1-\theta \mbox{ or } \theta \times Poisson(0\mid\lambda)\end{array} \right.
\end{equation}

% このデータが0のときに生成メカニズムが条件分岐していますから，ここに\verb|log_sum_exp|関数を入れる必要があります。
The generation mechanism branches when this data is zero, so you need to put the \verb|log_sum_exp| function here.
\begin{equation*}
   \verb|log_sum_exp| \begin{pmatrix} \log(1-\theta) \\ \log(\theta) + \verb|poisson_lpmf|(0\mid\lambda)\end{pmatrix}
\end{equation*}

% このことを念頭において，コードを書いてみましょう(コード\ref{stan::26_04})。
Keeping this in mind, let's try to write the code (code\ref{stan::26_04}).
\input{x26_mixed_escape_16}
% \paragraph{コード解説}
\paragraph{Code Explanation}
\begin{description}
	\item[dataブロック] データ長L，アウトカム変数Yを宣言しています。
	\item[parametersブロック] クローザーかどうかを決める$\theta$と，ポアソン分布のパラメータ$\lambda$がパラメータです。
	\item[modelブロック] データの中で，プログラミング言語としての\verb|if|文による条件分岐が出てきています\footnote{プログラミングでの条件分岐については第\ref{x16_Programming}講，セクション\ref{if_statement}，Pp.\pageref{if_statement}でやりましたね。忘れた人は戻って再確認しましょう。}。データが0であれば，クローザーでないか，ポアソン分布からゼロが出ているかです。そうでなければ，確率$\theta$経由でのポアソン分布です。
\end{description}


% これを実行するためのRコードはコード\ref{code:26_02}のようになります。
The R code to execute this would look like code \ref{code:26_02}.
\begin{lstlisting}[language=c++,label={code:26_02},caption={ゼロ過剰ポアソンモデルのコード}]
dat.tmp <- dat %>%
  dplyr::filter(position == "投手") %>%
  dplyr::filter(Games > 50) %>%
  dplyr::select(Save)

model <- cmdstan_model("ziPoisson.stan")
dataSet <- list(L = NROW(dat.tmp), Y = dat.tmp$Save)
fit <- model$sample(
  data = dataSet,
  chains = 4,
  parallel_chains = 4
)
\end{lstlisting}


% 推定した結果を見てみましょう。私の環境では出力\ref{MCMCout:out26_02}のようになりました。
Let's take a look at the estimated results. In my environment, it turned out like output \ref{MCMCout:out26_02}.
\begin{MCMCout}{}{out26_02}
	\begin{verbatim}
# A tibble: 2 × 7
  name         EAP       MED       MAP        SD       L95       U95
  <chr>  <num:.3!> <num:.3!> <num:.3!> <num:.3!> <num:.3!> <num:.3!>
1 lambda    14.666    14.665    14.661     0.261    14.154    15.185
2 theta      0.606     0.606     0.608     0.025     0.556     0.655
\end{verbatim}
\end{MCMCout}

% これを見ると，6割の確率でクローザー判定され，クローザーであれば$\lambda=14.7$ぐらいのポアソン分布でデータが生成されるということですね。
Upon observing this, there's a 60% chance that it will be determined as a 'closer'. If it is a 'closer', the data is generated according to a Poisson distribution with $\lambda$ approximately equal to 14.7.
% 事後予測分布のように，この確率モデルで出てくるデータを図\ref{fig:26_08}に示してみました。単一の分布では表現できないモデル化ができているのではないでしょうか。
As with the posterior predictive distribution, I tried to show the data that comes out in this probability model in Figure \ref{fig:26_08}. Aren't we able to model something that cannot be expressed in a single distribution?
\begin{figure}[ht]
   \begin{center}
      \includegraphics[width=0.8\linewidth,keepaspectratio]{../images/26_mixed/Rplot26_05.png}
      \caption{ゼロ過剰ポアソン分布による当てはめ。上がロウデータ，下が事後予測分布}
      \label{fig:26_08}
   \end{center}
\end{figure}

% もっとも，今回は1，2回セーブポイントがついている人のところをうまく表現できませんでした。しかしコードを書き換えることで，セーブ数0の先発投手と，本職のクローザーに加え，セーブ数が一桁ぐらいの中継ぎ投手の三種類に分割する，ということもできるでしょう。
However, this time we were unable to accurately represent those who have had a save point once or twice. However, by rewriting the code, it would be possible to divide into three types: starting pitchers with zero saves, professional closers, and relief pitchers with a single-digit number of saves.

% ここでみた分布の混ぜ方は，たとえば「お店に初来店した人とリピーターの来店数」とか，「20歳までの恋愛経験の数」のようなデータがあったときに，その特徴から分割してモデルを調整できます。研究のシーンでも，たとえば「質問紙で何も考えずにどちらとも言えないと答える人と，ちゃんと考えてどちらとも言えないに丸をつけた人」というように回答者のバイアスを除去してその特徴を考えるなど\footnote{この場合は中点の3が多い，というようなケースでしょうから，3過剰正規分布とでも言えば良いでしょうか。応用例として\citet{shimizu2018}などがあります。}，積極的な分析をできるのです。
The mixing of distributions we've seen here, for instance, allows us to adjust the model based on the characteristics when we have data such as "the number of visits from first-time customers and repeat customers to a store" or "the number of romantic relationships up to the age of 20". Even in the scene of research, for instance, it enables proactive analysis by considering characteristics after removing the bias of respondents such as "those who answer undecided without thinking in a questionnaire, and those who carefully think and mark undecided"\footnote{In this case, it would be a case where there are many midpoints of 3, so should we call it, an overabundance of 3 in a normal distribution? As an application example, there is \citet{shimizu2018}.}.


% さらに言えば，このポアソン分布の$\lambda$に線形モデルを入れて，ゼロ過剰ポアソン回帰分析にすることだってできますね。
Furthermore, speaking about this, it is possible to insert a linear model into the lambda of this Poisson distribution and make it a zero-inflated Poisson regression analysis.
% たとえばクローザーのセーブ数が年俸によって説明される，というモデルにするなら，コード\ref{stan::26_05}のようにすれば良いのです。
For example, if you model the number of saves by a closer being explained by their annual salary, then you should use the code as shown in \ref{stan::26_05}.
\begin{lstlisting}[language=C++,label={stan::26_05},caption={ゼロ過剰ポアソン回帰}]
data{
  int L;
  array[L] int Y;
  array[L] real X;
}

parameters{
  real<lower=0,upper=1> theta;
  real beta0;
  real beta1;
}

transformed parameters{
  real<lower=0> lambda[L];
  for(l in 1:L){
    lambda[l] = exp(beta0 + beta1 * X[l]);
  }
}

model{
  for(l in 1:L){
    if(Y[l]==0){
      target += log_sum_exp(log(1-theta),log(theta)+poisson_lpmf(0|lambda[l]));
    }else{
      target += log(theta) + poisson_lpmf(Y[l]|lambda[l]);
    }
  }
}
\end{lstlisting}

% これは一般化線形モデルの応用ですから，ポアソン分布に限らず使える技であることは，すぐにお分かりいただけることと思います。
This is an application of the generalized linear model, so you can easily understand that it is a technique that can be used not only for the Poisson distribution.

% データの分布や特徴を考えて，それにあった分析モデルを作っていくのは，モデラーとしての階段を一歩登ったことになります。図\ref{fig:26_09}には統計分析家の成長ステップの模式図を描いてみました\footnote{この図は東北大学大学院文学研究科の浜田宏先生が，専修大学社会知性開発研究センター主催の研究会(2019.2.26開催，於専修大学サテライトキャンパス)にてご講演いただいた際の，資料をもとに作図したものです。}。最初はデータの記述や線形モデルの当てはめなど，得られたデータに振り回されるような形でモデルを作り上げてきましたが，今回のように分布を混ぜ合わせると，さらにその表現力が豊かになることが実感できたのではないでしょうか。図にあるように，確率モデルのパラメータが心理事象の(数学的)記述から導出され，それがデータと合致する可動かで検証するようになれれば，さらにもう一段階研究のレベルが上がることになるでしょう。最終的には，心理学理論の方から新しい確率分布が導出される日が来るかもしれませんね。
Thinking about the distribution and characteristics of data, and creating an analysis model that suits them, means you have taken a step up the ladder as a modeler. In Figure \ref{fig:26_09}, I tried to draw a schematic diagram of the growth steps of a statistical analyst\footnote{This figure is based on the materials from a lecture given by Professor Hiroshi Hamada of the Graduate School of Arts and Letters, Tohoku University, at a research meeting sponsored by the Social Intelligence Development Research Center of Senshu University (held on February 26, 2019, at the Senshu University Satellite Campus)}. At first, we created models in a way that was influenced by the data we obtained, such as describing the data and fitting linear models. However, by mixing distributions as we did this time, you may have been able to feel that the expressiveness becomes richer. As shown in the figure, if the parameters of the probability model are derived from the (mathematical) description of psychological events and it is confirmed that they match the data, the level of research will rise another step. Eventually, a new probability distribution may be derived from psychological theories.
\begin{figure}[ht]
	\begin{center}
		\includegraphics[width=0.8\linewidth,keepaspectratio]{../images/26_mixed/slide26_04.png}
		\caption{統計モデラーの道。浜田宏先生(東北大学)のアイデアをもとにイラスト化したもの}
		\label{fig:26_09}
	\end{center}
\end{figure}

% \section{課題}以下のモデルを分析するR/Stanコードを提出してください。結果の解釈などを，スクリプトのコメントアウトや別添ファイルなどで提供してもらえるとなお良いです。もちろんRmdファイルでの提出であれば完璧です。なお提出されたコード単体でバグがなく動くことが確認できないものは，未提出扱いになります。コードの書き方などわからないところがあれば，曜日別TAか小杉までメールで連絡し，指導を受けてください。
\section{Task} Please submit the R/Stan code that analyzes the following model. It would be even better if you could provide interpretations of the results, such as commenting out the script or providing separate files. Of course, it’s perfect if you submit in Rmd file format. Please note that if we cannot confirm that the submitted code works without bugs on its own, it will be treated as if it was not submitted. If you have any questions about how to write the code, please email us on the TA or Kosugi by weekday, and receive guidance.

% \paragraph{混合分布モデル}今回は阪神Tigersのデータセットで行ったが，他の球団ではどうだろうか。別の球団データの年俸(対数をとったもの)のヒストグラムを描き，混合分布モデルが適用できそうなものを見つけたら，実際に当てはめて推定してみましょう。
\paragraph{Mixed Distribution Model} This time, we applied it to the data set of the Hanshin Tigers, but how would it be with other teams? Let's draw a histogram of the annual salary (logarithm) of other team data, and if you can find something that seems to apply a mixed distribution model, let's actually apply and estimate it.
% \paragraph{ゼロ過剰ポアソン回帰モデル}セーブポイントが年俸で説明できるとしたゼロ過剰ポアソン回帰を実行してみましょう。年俸データは標準化して単位を整えておくと良いでしょう。
\paragraph{Zero-Inflated Poisson Regression Model} Let's try running a zero-inflated Poisson regression, assuming that the save point can be explained by annual salary. It is a good idea to standardize your annual salary data to tidy up the units.
\end{document}
